% Web source: docs/05statistical_mechanics/02/02_010.md
\addcontentsline{toc}{section}{2-1　等重率の原理}

\begin{mondai}{2-1}{等重率の原理}{基本}
エネルギー $E$，体積 $V$，粒子数 $N$ が一定に保たれた孤立系を考える。これらの巨視的条件を満たす微視的状態の集合をミクロカノニカルアンサンブルとする。許されたすべての微視的状態は等しい確率で実現すると仮定し，以下の問いに答えよ。

\begin{enumerate}
  \item 系全体がとりうる微視的状態の総数を $\Omega$ とする。ある特定の状態 $k$ が実現する確率 $P_k$ を $\Omega$ を用いて表せ。また，すべての状態について確率を足し合わせ，規格化条件を満たすことを示せ。

  \item 各微視的状態 $k$ において，ある物理量 $A$ が値 $A_k$ をとるものとする。物理量 $A$ のアンサンブル平均 $\langle A \rangle$ と分散 $\langle A^2\rangle-\langle A\rangle^2$ を，$A_k$ と $\Omega$ を用いて表せ。
\end{enumerate}
\end{mondai}

\solutionheading
\begin{solutionparts}

\solutionpart{(1)}
許された微視的状態の総数が $\Omega$ であり，等重率の原理からすべての状態が等確率で実現するため，個々の特定の状態 $k$ が実現する確率 $P_k$ は状態数の逆数として一意に決まります。
\begin{equation*}
P_k = \frac{1}{\Omega}
\end{equation*}
すべての許された状態にわたって確率の総和を計算すると，以下のようになります。
\begin{equation*}
\sum_{k=1}^{\Omega} P_k = \sum_{k=1}^{\Omega} \frac{1}{\Omega} = \frac{1}{\Omega} \times \Omega = 1
\end{equation*}
これにより，全確率の総和が $1$ になるという規格化条件が満たされていることが示されました。

\solutionpart{(2)}
物理量 $A$ のアンサンブル平均 $\langle A \rangle$ は，各状態における値 $A_k$ に確率 $P_k = 1/\Omega$ を掛け，すべての微視的状態について足し合わせて求めます。
\begin{equation*}
\langle A \rangle = \sum_{k=1}^{\Omega} P_k A_k = \frac{1}{\Omega} \sum_{k=1}^{\Omega} A_k
\end{equation*}
同様に，$A^2$ の平均は
\begin{equation*}
\langle A^2\rangle=\frac{1}{\Omega}\sum_{k=1}^{\Omega}A_k^2
\end{equation*}
です。したがって，分散は
\begin{equation*}
\langle A^2\rangle-\langle A\rangle^2
=\frac{1}{\Omega}\sum_{k=1}^{\Omega}A_k^2
-\left(\frac{1}{\Omega}\sum_{k=1}^{\Omega}A_k\right)^2
\end{equation*}
となります。

\end{solutionparts}
